\section{Voice Agent 工程：低时延、多干扰、强鲁棒}

讲者明确表示，在当前阶段，可靠大规模语音 Agent 的主流形态仍是 pipeline：Speech-to-Text（ASR）$\rightarrow$ reasoning/orchestration $\rightarrow$ Text-to-Speech（TTS），而不是单一端到端 voice-to-voice。

\subsection{延迟预算与并行推理}

语音体验里最关键的是“用户停顿到 Agent 开口”的延迟。讲者强调要压缩每个 10ms 级别的开销，包括 endpoint detection、推理触发、TTS 首帧输出。为降低 P90/P95 波动，实践中会并发触发多个推理请求并采用先返回策略（hedging），再在后台取消慢请求。

\begin{knowledgebox}{语音响应延迟的简化分解}
\[
L_{\text{turn}} = L_{\text{EOU-detect}} + L_{\text{ASR}} + L_{\text{reason}} + L_{\text{TTS-first-byte}}
\]
其中真正影响主观体验的是 $L_{\text{turn}}$ 的高分位（如 P90/P95），而非均值。均值好看但尾部抖动严重，用户仍会感到“卡顿”和“抢话失败”。
\end{knowledgebox}

\subsection{打断检测与会话节奏控制}

人类对话中的“嗯嗯”“ok”常是 backchannel，不代表希望打断。系统若把所有声学事件都当中断，会出现频繁抢停、语义残缺。讲者提到需要专门训练中断检测模型，区分 acknowledgment 与真正的话轮夺取。

\begin{importantbox}{Voice UX 的关键不是“会说”，而是“会接话”}
高质量 voice agent 需要同时处理：端点检测、打断分类、重叠语音、多说话人、背景噪声、口音变化。任何一个环节失稳，都会放大成“你根本没听懂我”的体验崩塌。
\end{importantbox}

\subsection{ASR/TTS 评测与定制指标}

讲者指出 word error rate（WER）并不总是最佳指标。真实电话场景中，背景电视声、旁人说话、杂音都可能被“正确地不转写”。因此需要任务定制指标，评估“是否抓住主说话人关键信息并支撑正确动作”。

\begin{warningbox}{指标误配风险}
如果团队只优化通用 ASR 指标，可能得到“文字更完整、业务更失败”的反常结果。客服 Agent 必须按任务目标定义评测：身份校验字段、订单号关键槽位、动作触发准确率等。
\end{warningbox}

\subsection{本章小结}

语音 Agent 的难点不在单模型参数，而在端到端时序系统。要做到“听清、想清、说清”，必须用工程手段管理延迟尾部、话轮交互和多源噪声。

